% Generated draft: requires mathematical and source review.
\documentclass[11pt]{article}
\usepackage[margin=25mm]{geometry}
\usepackage{fontspec}
\setmainfont{Latin Modern Roman}
\usepackage{amsmath,amssymb,mathtools,unicode-math}
\setmathfont{Latin Modern Math}
\usepackage{enumitem,needspace,xurl,hyperref}
\setlength{\parindent}{0pt}
\setlength{\parskip}{5pt}
\setlength{\emergencystretch}{3em}
\newcommand{\N}{\mathbb{N}}
\newcommand{\Z}{\mathbb{Z}}
\newcommand{\Q}{\mathbb{Q}}
\newcommand{\R}{\mathbb{R}}
\newcommand{\C}{\mathbb{C}}
\newcommand{\F}{\mathbb{F}}
\newcommand{\A}{\mathbb{A}}
\newcommand{\PP}{\mathbb{P}}
\newcommand{\E}{\mathbb{E}}
\newcommand{\eps}{\varepsilon}
\newcommand{\dd}{\,\mathrm{d}}
\begin{document}
\section*{UnsolvedMath ID: 30006096}
\textbf{Source problem code:} catalogue entry
\subsection{Definitions and mathematical statement}
Let \(\mathbb S\) denote the sphere spectrum. For a space \(X\), let \(X_+\) be \(X\) with a disjoint basepoint and write
\[
D_+(X)=F(\Sigma^\infty X_+,\mathbb S)
\]
for its spherical cochain ring spectrum. Here \(F\) denotes the function spectrum; the supplied passage also uses the notation \(F(X_+,\mathbb S)\) and \(\mathbb S^X\). Let \(B\Z\) be the classifying space of the discrete group \(\Z\). The supplied setup is
\[
A=D_+(S^1)\simeq D_+(B\Z)\simeq\mathbb S^{h\Z},
\qquad R=\operatorname{THH}(A),
\]
where \(\Z\) acts trivially on \(\mathbb S\), the superscript \(h\Z\) denotes homotopy fixed points, and \(\operatorname{THH}\) denotes topological Hochschild homology. The equivalences use \(S^1\simeq B\Z\).

The passage states that for any \(\Z\)-equivariant \(A_\infty\) ring spectrum \(B\), there is a natural map
\[
R\longrightarrow\operatorname{THH}(B^{h\Z}).
\]
This is retained as a source assertion describing the setup. The precise conventions for equivariant \(A_\infty\) ring spectra, including units and compatibility of the action with multiplication, are not specified in the extract. The passage also describes the usual circle action on \(R\), calls this the \emph{outer circle}, and denotes it by \(T\). Its next sentence is truncated before explaining what this circle is to be distinguished from.

The passage attributes a role in the disproof of the Telescope Conjecture to \(R\), or its \(p\)-adic completion. This contextual claim is supplied attribution, not an independently verified conclusion here; neither a prime \(p\) nor a completion convention is specified.

\textbf{Unresolved mathematical target.} The supplied record contains definitions and descriptive assertions but no standalone question, conjecture, or requested conclusion. No precise open-problem statement can therefore be recovered without additional source material. The registry labels the record \texttt{open}, while its statement status is \texttt{unrecoverable}; its supplied verification note says that the report discusses structural results and does not pose a standalone open problem. These are supplied curation labels and assessments, not independently verified mathematical facts.
\subsection{Short English statement}
The extract, attributed to Neil P. Strickland, describes the topological Hochschild homology of spherical cochains on the circle and its outer circle action. It supplies no complete mathematical question, so this draft preserves the setup and explicitly leaves the intended problem unresolved.
\subsection{Sources}
\begin{itemize}
\item Supplied local UnsolvedMath registry record, permanent ID \(30006096\), titled \emph{Reported Structure of Circle-Cochain THH}; exact source problem code: \texttt{catalogue entry}. The record belongs to the Oberwolfach Reports Open Problems collection and gives proposed year \(2024\).
\item Supplied passage: \emph{THH of cochains on the circle}, attributed in the extract to Neil P. Strickland. Supplied report citation: \emph{Arbeitsgemeinschaft: Algebraic \(K\)-Theory and the Telescope Conjecture} (2025), Oberwolfach Reports, DOI 10.4171/owr/2024/47, pp.~2751--2804. Supplied source link: \url{https://doi.org/10.4171/owr/2024/47}. No passage-specific page number is supplied.
\item Supplied literature reference: \emph{Oberwolfach Report 47}, \url{https://doi.org/10.4171/OWR/2024/47}. The exact supplied literature-check date is \texttt{2026-08-22}. This is supplied dated curation, not an independently verified assessment. It labels the record open but states that the extract is a title and descriptive THH setup rather than a complete mathematical question or falsifiable claim, and consequently leaves its literature status unresolved.
\item The supplied statement-verification note says that the official report discusses structural results rather than a standalone open problem, that no separate clean statement is promoted, and that the original record is retained publicly for provenance. The background describes the literature triage as derived from supplied OWR curation and calls for independent verification of time-sensitive status and cited claims. No source was independently consulted for this draft.
\end{itemize}
\end{document}
